% The manuscript; stable-expansion.pdf beside it is this file built with pdflatex, run three times
% (sh tools/paper-build.sh stable-expansion in the GENChase repository).
% Status: DRAFT (2026-09-25). Not peer reviewed.
\documentclass[11pt]{article}
\usepackage[letterpaper,margin=1in]{geometry}
\usepackage[T1]{fontenc}
\usepackage{lmodern}
\usepackage{amsmath,amssymb,amsthm}
\usepackage{graphicx}
\usepackage{array}
\usepackage{booktabs}
\usepackage{microtype}
\usepackage[hidelinks]{hyperref}

\hypersetup{
  pdftitle={Stable Self-Similar Expansion of Four and Five Point Vortices and Confinement of Vortex Patches},
  pdfauthor={Chase Hendrick}}

\newtheorem{theorem}{Theorem}
\newtheorem{lemma}{Lemma}
\newtheorem{corollary}{Corollary}
\newtheorem{proposition}{Proposition}
\theoremstyle{definition}
\newtheorem{remark}{Remark}

\setlength{\emergencystretch}{2em}

\newcommand{\re}{\operatorname{Re}}
\newcommand{\im}{\operatorname{Im}}
\newcommand{\tr}{\operatorname{tr}}
\newcommand{\C}{\mathbb{C}}
\newcommand{\R}{\mathbb{R}}

\title{Stable Self-Similar Expansion of Four and Five Point Vortices and Confinement of Vortex Patches}
\author{Chase Hendrick\\[0.2em] {\small Independent Researcher}\\ {\small\href{mailto:chasewhendrick@gmail.com}{\texttt{chasewhendrick@gmail.com}}}\\ {\small ORCID \href{https://orcid.org/0009-0002-9754-6087}{0009-0002-9754-6087}}}
\date{\small Draft of 25 September 2026. Not peer reviewed.}

\begin{document}
\maketitle

\begin{abstract}
Point vortices can move self-similarly, keeping their shape while the configuration grows like $\sqrt{t}$ and turns. For three vortices such expanding configurations are stable, and Zbarsky used this to show that vortex patches placed at the three vortices stay confined for all time; he wrote that the result would most likely carry over to stably growing systems of four or more vortices, given sufficiently good stability. We give such configurations of four and five vortices, with the circulations $(-1, -5/2, -1/9, 4/5)$ and $(-1, 3/7, 7/8, -9/7, -47/35)$. In similarity variables their linearization has, besides the double eigenvalue $0$ of the rotation and of the one-parameter family of such configurations, only eigenvalues with real part $-1$ or $-2$. The proof is computer-assisted: the configurations are enclosed by the Krawczyk method in ball arithmetic, and the spectrum is controlled through an exact lemma, which shows that six eigenvalues are forced by the rotation, the scaling, the translations and the family and that the others come in pairs $k$, $2 - k$, together with enclosures of traces of powers of the Jacobian. The stability is nonlinear: every motion with the same circulations that starts near one of the two configurations stays within a bounded distance of an exactly self-similar expansion of a nearby member of the family, the member with the same energy, and an a priori estimate of the same kind holds for approximate solutions whose error is small and decays faster than the velocities. With this estimate in place of the one step of Zbarsky's proof that needs three vortices, his confinement theorem holds for the two configurations: vortex patches placed at the vortices stay within distance $t^{1/4 + \varepsilon}$ of their centres of vorticity for all time, and the centres move asymptotically like the self-similar expansion of a single member of the family. Linear stability is an open condition; in a random sample about one four-vortex collapse in seven and one five-vortex collapse in seventeen reverses into a linearly stable expansion.

\medskip
\noindent\textbf{Keywords:} point vortices; self-similar motion; nonlinear stability; vortex patches; computer-assisted proof; interval arithmetic.\\
\textbf{MSC 2020:} 76B47, 35Q31, 37N10, 65G20.
\end{abstract}

\section{Introduction}

$N$ point vortices with circulations $\Gamma_j$ at positions $z_j \in \C$ move by
\begin{equation}\label{eq:bs}
  \overline{\dot z_j} = \frac{1}{2\pi i}\sum_{k \ne j}\frac{\Gamma_k}{z_j - z_k}.
\end{equation}
A motion is self-similar when the configuration keeps its shape while it grows or shrinks and turns; this happens when $\dot z_j = \kappa(z_j - z_c)$ for all $j$ at $t = 0$, with one $\kappa \in \C$, where $z_c = \sum_j\Gamma_jz_j/\sum_j\Gamma_j$ \cite[Sect.~2]{mw2026}. Then the size of the configuration grows or shrinks like $|1 + 2t\re\kappa|^{1/2}$, the configuration expands when $\re\kappa > 0$ and collapses in finite time when $\re\kappa < 0$ \cite{kimura1987, ns1979, aref2010, gotoda2021, ks2022}. Reversing every circulation reverses time, so each collapse becomes an expansion of the same shape, and conversely.

Zbarsky \cite{zbarsky2021} proved that for a generic self-similarly expanding configuration of three vortices one may replace each vortex by a small patch of vorticity and the patches stay confined for all time, their radius growing more slowly than their distance. The restriction to three vortices comes from the stability analysis. He writes that ``if one found stably growing systems of 4 or more vortices, the confinement result would most likely carry over with little modification. However, one needs sufficiently good stability results; orbital linear stability, which may not be hard to obtain for some systems, is not enough'' \cite[pp.~4--5 of the arXiv version]{zbarsky2021}; the same point recurs in \cite[after Cor.~16]{zbarsky2024}. For three vortices, stability of the expanding family goes back to Tavantzis and Ting \cite{tt1988}; Iwayama and Yajima studied the linear stability of the self-similar motions of three vortices in the $\alpha$-models \cite{iy2023}.

We show that stable expansions of four and five vortices exist, and that Zbarsky's confinement theorem holds for them. Section~\ref{sec:setting} sets up similarity variables, in which a self-similar motion is a fixed point, and proves a lemma on the structure of the spectrum of the linearization (Lemma~\ref{lem:spectrum}): six eigenvalues are forced by the rotation, the scaling, the translations and the one-parameter family of self-similar configurations with the given circulations (the second eigenvalue $2$ through a pairing), and the others come in pairs $k$, $2 - k$. For three vortices nothing else is left. For four vortices one pair is left, and a single number decides its stability; for five, two pairs. Section~\ref{sec:results} states and proves, with a computer-assisted proof, that there are configurations of four and of five vortices whose reversal expands with every other exponent of real part $1$ (Theorems~\ref{thm:four} and~\ref{thm:five}, Corollary~\ref{cor:expansion}). Section~\ref{sec:nonlinear} proves that the stability is nonlinear, with the sharp rate and a limit fixed by the energy (Theorem~\ref{thm:nonlinear}, Corollary~\ref{cor:physical}), together with an a priori estimate for approximate solutions (Proposition~\ref{prop:forced}). Section~\ref{sec:patches} uses that estimate in place of the one step of Zbarsky's proof that needs three vortices and obtains confinement of vortex patches for the two configurations (Theorem~\ref{thm:patches}). Section~\ref{sec:numerics} gives numerical illustrations, controls and a sample of how common such configurations are, and Section~\ref{sec:discussion} discusses what is and is not shown.

\section{Similarity variables and the structure of the spectrum}\label{sec:setting}

\subsection{Similarity variables}

Let $V(\zeta) = (V_1, \ldots, V_N)$ be the velocity given by \eqref{eq:bs}, $V_j = (i/2\pi)\sum_{k \ne j}\Gamma_k/\overline{(\zeta_j - \zeta_k)}$. It is homogeneous of degree $-1$, $V(\lambda\zeta) = \lambda^{-1}V(\zeta)$ for real $\lambda > 0$, and equivariant under rotations, $V(e^{i\phi}\zeta) = e^{i\phi}V(\zeta)$. For $\kappa = a + i\Omega$ put $z_j(t) = \rho(t)e^{i\Omega\tau}\zeta_j(\tau)$ with $\rho^2 = 1 + 2at$ and $d\tau/dt = \rho^{-2}$; then \eqref{eq:bs} becomes
\begin{equation}\label{eq:sim}
  \frac{d\zeta}{d\tau} = V(\zeta) - \kappa\zeta ,
\end{equation}
and a self-similar motion with this $\kappa$ and $z_c = 0$ is a fixed point of \eqref{eq:sim}. We fix the scale of time by $2\pi a = -1$ for a collapse, write $2\pi\kappa = -1 + ib$, and work with $E(\zeta) = 2\pi(V(\zeta) - \kappa\zeta) = 2\pi V(\zeta) + (1 - ib)\zeta$ and the rescaled time $s = \tau/2\pi$; for a collapse, $s = -\ln\rho$. The winding number of the collapse is $P = |b|/2$ \cite{mw2026}. Let $DE$ be the real $2N \times 2N$ Jacobian of $E$ at a zero $\zeta^*$. A perturbation along an eigenvector with eigenvalue $k$ behaves like $e^{ks} = \rho^{-k}$ relative to the size $\rho$ of the configuration. Writing $\delta\zeta = u + iv$, one finds
\begin{equation}\label{eq:DE}
  DE = \begin{pmatrix} A_r + I & A_i + bI\\ A_i - bI & -A_r + I\end{pmatrix},\qquad A_{jk} = \frac{i\Gamma_k}{\overline{(\zeta_j - \zeta_k)}^{\,2}}\ (k \ne j),\quad A_{jj} = -\sum_{k \ne j}A_{jk},
\end{equation}
with $A = A_r + iA_i$ evaluated at $\zeta^*$.

\emph{Time reversal.} If $\zeta^*$ is a collapse for the circulations $\Gamma$, it is a fixed point of the similarity dynamics of the circulations $-\Gamma$ with $\kappa$ replaced by $-\kappa$, an expansion with the same winding, and the Jacobian of its similarity dynamics is $-DE$; the similarity time is then $s = \ln\rho$. So a mode with eigenvalue $k$ of $DE$ behaves in the expansion like $\rho^{-k}$ relative to its size, and it decays as the configuration grows exactly when $\re k > 0$.

\subsection{The structure of the spectrum}

\begin{lemma}\label{lem:spectrum}
Let $\zeta^*$ be a zero of $E$, with all $\Gamma_j \ne 0$ and the $\zeta_j^*$ distinct, and suppose that near $(\zeta^*, b)$ the zeros $(\zeta, b)$ of $E$, with the circulations fixed, form up to rotation a smooth curve through $(\zeta^*, b)$ whose $\zeta$-component is transverse to the rotations. Then the characteristic polynomial $p$ of $DE$ satisfies $p(\lambda) = p(2 - \lambda)$ and is divisible by
\[
  \lambda^2(\lambda - 2)^2\bigl((\lambda - 1)^2 + b^2\bigr),
\]
and the quotient $q$, of degree $2N - 6$, also satisfies $q(\lambda) = q(2 - \lambda)$. In particular the $2N - 6$ remaining eigenvalues, counted with multiplicity, come in pairs $k$, $2 - k$.
\end{lemma}

\begin{proof}
\emph{Rotation.} $E(e^{i\phi}\zeta) = e^{i\phi}E(\zeta)$, so differentiating at $\phi = 0$ gives $DE(i\zeta^*) = iE(\zeta^*) = 0$.

\emph{Scaling.} $E(\lambda\zeta) = 2\pi V(\zeta)/\lambda + (1 - ib)\lambda\zeta$; at $\lambda = 1$ the derivative is $-2\pi V(\zeta^*) + (1 - ib)\zeta^* = 2(1 - ib)\zeta^*$, since $2\pi V(\zeta^*) = -(1 - ib)\zeta^*$. So $DE\,\zeta^* = 2\zeta^* - 2b\,i\zeta^*$, and on the span of $\zeta^*$ and $i\zeta^*$ the map $DE$ is triangular with eigenvalues $2$ and $0$.

\emph{Translations.} $V$ is invariant under $\zeta \mapsto \zeta + c(1, \ldots, 1)$, so $DE$ acts on $\mathbf 1 = (1, \ldots, 1)$ and $i\mathbf 1$ as multiplication by $1 - ib$: $DE\,\mathbf 1 = \mathbf 1 - b\,i\mathbf 1$, $DE(i\mathbf 1) = i\mathbf 1 + b\mathbf 1$, with eigenvalues $1 \pm ib$.

\emph{The family.} Along the curve of zeros, $E(\xi(\sigma); b(\sigma)) = 0$, and $\partial_bE = -i\zeta$. Differentiating, $DE\,\xi'(0) = b'(0)\,i\zeta^*$, so $DE^2\xi'(0) = 0$. As $\xi'(0)$ is not a real multiple of $i\zeta^*$, the generalized kernel of $DE$ has dimension at least $2$: the algebraic multiplicity of $0$ is at least $2$.

\emph{Pairing.} With $H = -\frac{1}{2}\sum_{j \ne k}\Gamma_j\Gamma_k\ln|\zeta_j - \zeta_k|$ and $I = \sum_j\Gamma_j|\zeta_j|^2$, the field $\zeta \mapsto 2\pi V(\zeta) - ib\zeta$ is the Hamiltonian vector field of $H + (b/2)I$ for the symplectic form $\sum_j\Gamma_j\,dx_j \wedge dy_j$, which is nondegenerate because every $\Gamma_j \ne 0$ (Kirchhoff's form of \eqref{eq:bs} and the rotations generated by $I$, up to the sign convention of the form). The linearization of a Hamiltonian vector field at any point is an infinitesimally symplectic matrix, whose characteristic polynomial is even. Since $DE$ is this linearization plus the identity, $p(\lambda) = r(\lambda - 1)$ with $r$ even, that is $p(\lambda) = p(2 - \lambda)$, with multiplicities.

\emph{Conclusion.} By pairing, the multiplicity of $2$ equals that of $0$, at least $2$, so $\lambda^2(\lambda - 2)^2$ divides $p$. The span of $\mathbf 1$ and $i\mathbf 1$ is invariant, and the characteristic polynomial of $DE$ on it, $(\lambda - 1)^2 + b^2$, divides $p$ too; its roots $1 \pm ib$ are never $0$ or $2$, so the two factors are coprime and $p$ is divisible by the displayed product. That product is invariant under $\lambda \mapsto 2 - \lambda$, hence so is the quotient.
\end{proof}

For three vortices $2N - 6 = 0$: in similarity variables every exponent of a three-vortex expansion is forced by the symmetries and the family, which is the linear shadow of the classical stability of the expanding three-vortex family \cite{tt1988}. For four vortices one pair $k$, $2 - k$ remains, and with $c = k(2 - k)$ the sum of squares of all eigenvalues gives
\begin{equation}\label{eq:c}
  \tr(DE^2) = \bigl(0 + 0 + 4 + 4 + 2(1 - b^2)\bigr) + (4 - 2c),\qquad c = 3 + 3b^2 - \tr(A\bar A),
\end{equation}
using $\tr(DE^2) = 2\tr(A\bar A) + 2N(1 - b^2)$, which follows from \eqref{eq:DE}. Both remaining exponents have positive real part exactly when $c > 0$; when $c > 1$ they are $1 \pm i\sqrt{c - 1}$. For five vortices two pairs remain; with $u = (k - 1)^2$ they are the roots $u_1$, $u_2$ of $u^2 - m_1u + m_2$, where $m_1 = u_1 + u_2$ and $u_1^2 + u_2^2$ follow from $\tr((DE - I)^2)$ and $\tr((DE - I)^4)$ after the contributions $4 - 2b^2$ and $4 + 2b^4$ of the six forced eigenvalues are removed. Every remaining exponent has positive real part exactly when every root $u$ satisfies $|\re\sqrt{u}| < 1$; in particular it suffices that both roots are real and less than $1$.

\section{Results}\label{sec:results}

\begin{theorem}[Four vortices]\label{thm:four}
For the circulations $\Gamma = (1, 5/2, 1/9, -4/5)$ there is a self-similar collapse of four Euler point vortices with $2\pi\kappa = -1 + ib$, winding
\[
  P = b/2 = 2.5258464741344909107270\ldots,
\]
and $z_4$ with real part $8/25$ in the gauge $z_1 \in \R$, $z_c = 0$, whose Jacobian $DE$ has the six eigenvalues $0$, $0$, $2$, $2$, $1 \pm ib$ of Lemma~\ref{lem:spectrum} and the pair $1 \pm i\vartheta$ with $\vartheta = 4.4588620962481\ldots$; here $c = 20.881451193358\ldots$ in \eqref{eq:c}.
\end{theorem}

\begin{theorem}[Five vortices]\label{thm:five}
For the circulations $\Gamma = (1, -3/7, -7/8, 9/7, 47/35)$ there is a self-similar collapse of five Euler point vortices with winding $P = 2.40234703541713841\ldots$ and $z_2$ with real part $3/5$ in the same gauge, whose Jacobian $DE$ has the six eigenvalues of Lemma~\ref{lem:spectrum} and the four eigenvalues $1 \pm 3.4538555\ldots i$ and $1 \pm 3.1885675\ldots i$.
\end{theorem}

\begin{corollary}\label{cor:expansion}
For the circulations $(-1, -5/2, -1/9, 4/5)$ and $(-1, 3/7, 7/8, -9/7, -47/35)$ the configurations of Theorems~\ref{thm:four} and~\ref{thm:five} expand self-similarly, growing like $\sqrt{t}$. In similarity variables the linearization at them has the double eigenvalue $0$, whose generalized eigenspace consists of the rotation and the tangent of the one-parameter family of expanding configurations with the same circulations, and every other eigenvalue has real part $-1$ or $-2$. In the invariant complement of the rotation and the family, every solution of the linearized equation therefore decays, relative to the size $\rho$ of the configuration, at least like $\rho^{-1}$, that is like $t^{-1/2}$.
\end{corollary}

\begin{proof}[Proof of the theorems]
The program \texttt{verify\_\allowbreak stable\_\allowbreak expansion.py} (Section~\ref{sec:numerics}) works in the gauge of \cite[Sect.~8]{mw2026}: $\Gamma_1 = 1$, $z_1$ real, $\sum_k\Gamma_k/(z_j - z_k) = \lambda\overline{z_j}$ with $\lambda = 2P - i$, which is $E = 0$ with $b = 2P$ and $z_c = 0$. For four vortices it fixes $x_4 = 8/25$, $\Gamma_2 = 5/2$ and $\Gamma_3 = 1/9$ and proves by the Krawczyk test, in the ball arithmetic of FLINT/Arb, that the $8$ equations have exactly one zero in a box of radius $10^{-4}$ around the numerical solution; on the box the circulations are nonzero, the total circulation is nonzero, the vortices are distinct and away from the origin, and the necessary conditions $\sum_j\Gamma_jz_j = 0$, $\sum_{j<k}\Gamma_j\Gamma_k = 0$ and $\sum_j\Gamma_j|z_j|^2 = 0$ are enclosed with $0$. The second of these forces $\Gamma_4 = -4/5$. The Jacobian of the square system is invertible on the box, so by the implicit function theorem the zeros with $\Gamma_2$, $\Gamma_3$ fixed form a smooth curve parametrized by $x_4$; along it $\Gamma_4$ stays $-4/5$, and its tangent $v$ has $\im v_1 = 0$ and $v_4$ with real part $1$, while $i\zeta^*$ has $\im(i\zeta_1^*) = x_1 \ne 0$, so $v$ is not a real multiple of $i\zeta^*$. The hypotheses of Lemma~\ref{lem:spectrum} hold, and the program also checks on the box that the vectors of its proof satisfy their identities. It then encloses $c = 3 + 3b^2 - \tr(A\bar A)$ by $20.8814511933587 \pm 1.4\cdot10^{-14}$, so $c > 1$ and the remaining pair is $1 \pm i\sqrt{c - 1}$; a second enclosure from $\tr((DE - I)^2)$ agrees. For five vortices it fixes $x_2 = 3/5$ and $\Gamma_2$, $\Gamma_3$, $\Gamma_4$, proves existence the same way, finds $\Gamma_5 = 47/35$, checks the hypotheses of the lemma, and encloses $m_1$ and $m_2$; the discriminant $m_1^2 - 4m_2$ is positive and both roots, $u_1 = -11.929118\ldots$ and $u_2 = -10.166963\ldots$, are negative, so the four remaining eigenvalues are $1 \pm i\sqrt{-u_1}$ and $1 \pm i\sqrt{-u_2}$.
\end{proof}

\begin{proof}[Proof of the corollary]
By the time reversal of Section~\ref{sec:setting} the configurations expand for the negated circulations, with Jacobian $-DE$ in similarity variables and similarity time $s = \ln\rho$. The eigenvalues of $-DE$ are $0$, $0$, $-2$, $-2$, $-1 \mp ib$ and, by the theorems, $-1 \mp i\vartheta$ or $-1 \mp i\sqrt{-u_{1,2}}$. By the theorems no remaining eigenvalue is $0$, so the generalized eigenspace of $0$ is exactly two-dimensional and, by the proof of Lemma~\ref{lem:spectrum}, spanned by $i\zeta^*$ and the tangent of the family. The eigenvalues with real part $-1$ are simple: $\vartheta \ne b$ for four vortices and $\sqrt{-u_1}$, $\sqrt{-u_2}$ and $b$ are distinct for five, which the program checks. The eigenvalue $-2$ is double, the partner of the one at $0$, and even as a Jordan block it contributes only $s\,e^{-2s}$. On the complementary invariant subspace every solution of the linearized equation therefore decays at least like $e^{-s} = \rho^{-1}$, and $\rho^2$ grows linearly in $t$.
\end{proof}

\begin{remark}[Open sets]\label{rem:open}
The conclusions are open conditions: $c$, and the roots $u$, depend continuously on the configuration and the circulations along the collapse set, which near these points is a smooth manifold of dimension $N - 1$ (modulo translations, rotations and the scalings of lengths and circulations) \cite[Sect.~7]{mw2026}. So every collapse near those of the theorems, with nearby circulations, reverses into a linearly stable expansion as well; the stable expansions of four and five vortices form open subsets of the $(N - 1)$-dimensional collapse set, not isolated examples.
\end{remark}

\section{Nonlinear stability}\label{sec:nonlinear}

Throughout this section $\zeta^*$ is the configuration of Theorem~\ref{thm:four} or of Theorem~\ref{thm:five}, $b^*$ is its $b$, and the circulations are negated, so that $\zeta^*$ expands. By the time reversal of Section~\ref{sec:setting}, a motion $z(t)$ of the vortices with circulations $-\Gamma$ is $z(t) = \rho(t)e^{-ib^*s}\zeta(s)$ with $\rho^2 = 1 + t/\pi$, $s = \ln\rho$ and
\begin{equation}\label{eq:exp}
  \frac{d\zeta}{ds} = -E(\zeta),\qquad E(\zeta) = 2\pi V(\zeta) + (1 - ib^*)\zeta,
\end{equation}
where $V$ and $E$ are formed with the circulations $\Gamma$ of the theorems; this holds for every motion, not only for self-similar ones. We use the Euclidean norm $|\cdot|$ on $\C^N = \R^{2N}$, write $\mathcal O = \{e^{i\phi}\zeta^*: \phi \in \R\}$ for the orbit of $\zeta^*$ under rotations, and $H = -\sum_{j<k}\Gamma_j\Gamma_k\ln|\zeta_j - \zeta_k|$ for the energy, which is the same for $\Gamma$ and $-\Gamma$. By the proof of the theorems, the zeros of $E(\zeta; b) = 2\pi V(\zeta) + (1 - ib)\zeta$ near $(\zeta^*, b^*)$ with $z_1$ real form a real-analytic curve $\sigma \mapsto (\xi(\sigma), b(\sigma))$, $|\sigma| < \sigma_1$, with $\xi(0) = \zeta^*$ and $b(0) = b^*$, parametrized by $\sigma = x_4 - 8/25$ or $\sigma = x_2 - 3/5$. Since $E(\zeta; b^*) = E(\zeta; b) + i(b - b^*)\zeta$, the solution of \eqref{eq:exp} starting at $\xi(\sigma)$ is $e^{-i(b(\sigma) - b^*)s}\xi(\sigma)$: another self-similar expansion, with a different shape and winding. We write $\xi'$ for the derivative in $\sigma$ and $\zeta'$ for the derivative of a solution in $s$. We fix $\sigma_0 \in (0, \sigma_1)$ once, small enough for Lemma~\ref{lem:coords} below and such that $H(\xi(\sigma))$ is strictly monotone for $|\sigma| \le \sigma_0$; the latter is possible because the program encloses $dH(\xi(\sigma))/d\sigma$ at $\sigma = 0$ away from $0$ ($-0.5133915789\ldots$ for four vortices and $0.2322290983\ldots$ for five).

\begin{theorem}[Nonlinear stability]\label{thm:nonlinear}
There are $\delta_0 > 0$ and $C$ such that every solution of \eqref{eq:exp} with $\delta := \operatorname{dist}(\zeta(0), \mathcal O) \le \delta_0$ exists for all $s \ge 0$, and there are $\sigma_\infty$ with $|\sigma_\infty| \le C\delta$ and $\phi_\infty \in \R$ with
\[
  \bigl|\zeta(s) - e^{i(\phi_\infty - (b(\sigma_\infty) - b^*)s)}\xi(\sigma_\infty)\bigr| \le C\delta e^{-s}\qquad (s \ge 0).
\]
Moreover $\sigma_\infty$ is the unique $\sigma$ with $|\sigma| \le \sigma_0$ and $H(\xi(\sigma)) = H(\zeta(0))$.
\end{theorem}

The constants $\delta_0$ and $C$ are not explicit.

\begin{corollary}[In the original variables]\label{cor:physical}
Let $z(t)$ be a motion of $N$ point vortices with the circulations $-\Gamma$ and $|z(0) - e^{i\phi}\zeta^*| \le \delta \le \delta_0$ for some $\phi$. Then there are a member $\xi(\sigma_\infty)$ of the family, with the energy $H(z(0))$, and $\phi_\infty \in \R$ such that the exactly self-similar expansion $Z(t) = \rho(t)e^{i(\phi_\infty - b(\sigma_\infty)\ln\rho(t))}\xi(\sigma_\infty)$ satisfies $|z(t) - Z(t)| \le C\delta$ for all $t \ge 0$. The vortices stay within a bounded distance of an exactly self-similar expansion whose size grows like $\sqrt t$, and the relative deviation is $O(t^{-1/2})$.
\end{corollary}

\begin{proof}
Here $\zeta(0) = z(0)$. Multiplying the estimate of Theorem~\ref{thm:nonlinear} by $\rho e^{-ib^*s} = e^{s}e^{-ib^*s}$ gives $|z(t) - Z(t)|\le C\delta$. $Z$ is a motion because $\xi(\sigma_\infty)$ is a zero of $E(\cdot\,; b(\sigma_\infty))$, so a fixed point of the similarity dynamics with $b(\sigma_\infty)$ in place of $b^*$, and $H(\xi(\sigma_\infty)) = H(\zeta(0))$ by the theorem.
\end{proof}

The theorem follows from an a priori estimate for approximate solutions, whose constants do not depend on how long the approximate solution is followed; that is the form a continuation argument needs (Section~\ref{sec:discussion}).

\begin{proposition}[Forced version]\label{prop:forced}
For $\mu > 0$ and $0 < \nu < \min(\mu, 1)$ there are $\delta_1 > 0$ and $C_1$, depending only on $\zeta^*$, $\mu$ and $\nu$, with the following property. Let $0 < S \le \infty$ and let $\zeta: [0, S) \to \C^N$ be continuously differentiable with
\[
  \bigl|\zeta'(s) + E(\zeta(s))\bigr| \le \varepsilon e^{-\mu s}\quad(0 \le s < S),\qquad \delta := \operatorname{dist}(\zeta(0), \mathcal O),\qquad \delta + \varepsilon \le \delta_1.
\]
Then for $0 \le s < S$ we have $\operatorname{dist}(\zeta(s), \mathcal O) \le C_1(\delta + \varepsilon)$, and the distance from $\zeta(s)$ to $\{e^{i\phi}\xi(\sigma): \phi \in \R,\ |\sigma| \le \sigma_0\}$ is at most $C_1(\delta + \varepsilon)e^{-\nu s}$. If $S = \infty$, there are $\sigma_\infty$ with $|\sigma_\infty| \le C_1(\delta + \varepsilon)$ and a continuous function $\phi$ with $|\zeta(s) - e^{i\phi(s)}\xi(\sigma_\infty)| \le C_1(\delta + \varepsilon)e^{-\nu s}$.
\end{proposition}

Suppose that $z$ satisfies $\dot z = V_{-\Gamma}(z) + e(t)$ for $t \ge 0$, where $V_{-\Gamma}$ is the velocity \eqref{eq:bs} for the circulations $-\Gamma$ and $e$ is a continuous forcing with $|e(t)| \le K(1 + t)^{-(1+\mu)/2}$. Then $\zeta(s) = \rho^{-1}e^{ib^*s}z(t)$ satisfies $\zeta' = -E(\zeta) + 2\pi\rho e^{ib^*s}e(t)$, and since $1 + t \ge \rho^2$ the forcing is at most $2\pi Ke^{-\mu s}$: the proposition applies with $\varepsilon = 2\pi K$ when $\delta + 2\pi K \le \delta_1$. So the error must be small and decay faster than the velocities, which are of order $t^{-1/2}$.

That a manifold of equilibria whose transverse spectrum lies in the open left half-plane attracts every nearby solution exponentially, each to one of its equilibria, is classical \cite{aulbach1984}. We give the short proof because the family is a curve of equilibria only modulo the rotation, and because we need the forced version and the sharp rate.

\begin{lemma}[Coordinates]\label{lem:coords}
Let $W_0 = \operatorname{span}_\R\{i\zeta^*, \xi'(0)\}$ and let $W_1$ be the sum of the generalized eigenspaces of $DE$ at $\zeta^*$ for its nonzero eigenvalues, and fix a linear isomorphism $Q: \R^{2N-2} \to W_1$. For $\sigma_0$, $r_0 > 0$ small enough, $\Phi(\phi, \sigma, y) = e^{i\phi}(\xi(\sigma) + Qy)$ is a real-analytic diffeomorphism from $(\R/2\pi\mathbb Z) \times (-\sigma_0, \sigma_0) \times \{|y| < r_0\}$ onto an open neighborhood $U$ of $\mathcal O$. In these coordinates a continuously differentiable $\zeta(s) = \Phi(\phi, \sigma, y) \in U$ satisfies
\begin{equation}\label{eq:coords}
  (\phi', \sigma', y') = \bigl(h(\sigma, y), f(\sigma, y), g(\sigma, y)\bigr) + L_{\sigma, y}^{-1}\bigl(e^{-i\phi}(\zeta' + E(\zeta))\bigr),
\end{equation}
where $L_{\sigma,y}(p_1, p_2, w) = p_1\,i(\xi(\sigma) + Qy) + p_2\,\xi'(\sigma) + Qw$ and $(h, f, g)(\sigma, y) = -L_{\sigma,y}^{-1}E(\xi(\sigma) + Qy)$ are real-analytic, with $h(\sigma, 0) = -(b(\sigma) - b^*)$, $f(\sigma, 0) = 0$ and $g(\sigma, 0) = 0$. The spectrum of $B(\sigma) = \partial_yg(\sigma, 0)$ consists of the eigenvalue $-2$, of algebraic multiplicity $2$, and of $2N - 4$ simple eigenvalues on the line $\re\lambda = -1$, and there is $K_B$ with $\|e^{B(\sigma)s}\| \le K_Be^{-s}$ for $s \ge 0$ and $|\sigma| \le \sigma_0$.
\end{lemma}

\begin{proof}
By Corollary~\ref{cor:expansion} the generalized kernel of $DE$ is $W_0$, so $\C^N = W_0 \oplus W_1$ with both summands invariant under $DE$. The differential of $\Phi$ at $(0, 0, 0)$ is $L_{0,0}$, which is invertible; by rotation equivariance the differential of $\Phi$ is invertible at every $(\phi, \sigma, y)$ with $\sigma$, $y$ small. If $\Phi(\phi, \sigma, y) = \Phi(\tilde\phi, \tilde\sigma, \tilde y)$ with small $\sigma$, $\tilde\sigma$, $y$, $\tilde y$, then $e^{i(\tilde\phi - \phi)}\zeta^*$ is close to $\zeta^*$, so $\tilde\phi - \phi$ is close to $2\pi\mathbb Z$ because $\zeta^* \ne 0$, and the inverse function theorem at $(0, 0, 0)$ gives equality. This proves the first claim. Differentiating $\zeta = \Phi(\phi, \sigma, y)$ gives $\zeta' = e^{i\phi}L_{\sigma,y}(\phi', \sigma', y')$, and $E(\zeta) = e^{i\phi}E(\xi(\sigma) + Qy)$ by equivariance, which gives \eqref{eq:coords}. At $y = 0$, $E(\xi(\sigma)) = i(b(\sigma) - b^*)\xi(\sigma) = L_{\sigma, 0}(b(\sigma) - b^*, 0, 0)$, which gives the values of $h$, $f$ and $g$ there.

For the spectrum, consider \eqref{eq:coords} without forcing, with $b^*$ replaced by $b(\sigma)$ for a fixed $\sigma$. Rotation equivariance shows that this changes only the equation for $\phi$, by an additive constant, and $\xi(\sigma)$ becomes a fixed point. Its linearization in the coordinates is block upper triangular with diagonal blocks $\bigl(\begin{smallmatrix}0 & *\\ 0 & 0\end{smallmatrix}\bigr)$ and $B(\sigma)$, because $f(\cdot, 0) = g(\cdot, 0) = 0$, and it is similar to $-DE(\xi(\sigma); b(\sigma))$. Lemma~\ref{lem:spectrum} applies at every point of the curve, so the spectrum of $B(\sigma)$ consists of $-2$, $-2$, $-1 \pm ib(\sigma)$ and the negatives of the remaining eigenvalues. At $\sigma = 0$ these lie on $\re\lambda = -1$ and are simple, by the theorems. Along the curve they stay there for small $|\sigma|$: for four vortices $c > 1$ is preserved, and for five a positive discriminant and two negative roots $u$ are preserved, while being distinct is an open condition. The eigenvalue $-2$ keeps algebraic multiplicity $2$ because the others stay away from it. So
\[
  e^{B(\sigma)s} = \sum_{j=1}^{2N-4}e^{\lambda_j(\sigma)s}\Pi_j(\sigma) + e^{-2s}\bigl(\mathrm{Id} + sM(\sigma)\bigr)\Pi_0(\sigma),
\]
where the $\Pi_j$ are the spectral projections of the simple eigenvalues $\lambda_j$, $\re\lambda_j = -1$, $\Pi_0$ is that of $-2$ and $M = (B + 2)\Pi_0$ satisfies $M^2 = 0$. The projections are continuous in $\sigma$, so $K_B$ exists after shrinking $\sigma_0$.
\end{proof}

\begin{proof}[Proof of Proposition~\ref{prop:forced}]
Choose $\gamma \in (\nu, 1)$ with $\gamma \ne \mu$. The spectrum of $B(0)$ lies in $\re\lambda \le -1$, so there is an inner product on $\R^{2N-2}$, with norm $|\cdot|_*$, in which $\langle B(0)y, y\rangle_* \le -\frac{1 + \gamma}{2}|y|_*^2$. Since $B(\sigma) \to B(0)$ and $g(\sigma, y) = B(\sigma)y + O(|y|^2)$, there are $\sigma_2 \in (0, \sigma_0]$ and $r_2 > 0$, depending on $\gamma$, such that $\langle g(\sigma, y), y\rangle_* \le -\gamma|y|_*^2$, $|f(\sigma, y)| \le C_f|y|_*$ and $\|L_{\sigma,y}^{-1}\| \le C_L$ for $|\sigma| \le \sigma_2$ and $|y|_* \le r_2$, a part of the domain of Lemma~\ref{lem:coords}. By the inverse function theorem and equivariance, $\operatorname{dist}(\zeta(0), \mathcal O) = \delta$ gives $\zeta(0) = \Phi(\phi_0, \sigma(0), y(0))$ with $|\sigma(0)| + |y(0)|_* \le C_0\delta$ when $\delta$ is small.

Let $S_1 \le S$ be the supremum of the $s_1$ such that $\zeta(s) \in \Phi(\R \times [-\sigma_2/2, \sigma_2/2] \times \{|y|_* \le r_2/2\})$ for $0 \le s \le s_1$. On $[0, S_1)$ the forcing term in \eqref{eq:coords} has size at most $C_L\varepsilon e^{-\mu s}$ (in the norm $|\cdot|_*$, after enlarging $C_L$), so the upper right Dini derivative of $|y|_*$ satisfies $D^+|y|_* \le -\gamma|y|_* + C_L\varepsilon e^{-\mu s}$; where $y = 0$ this holds because $y'$ is then the forcing term alone. The comparison lemma gives
\[
  |y(s)|_* \le e^{-\gamma s}|y(0)|_* + C_L\varepsilon\int_0^se^{-\gamma(s - \tau)}e^{-\mu\tau}\,d\tau \le C_0\delta e^{-\gamma s} + \frac{C_L\varepsilon}{|\gamma - \mu|}e^{-\min(\gamma, \mu)s} \le C_2(\delta + \varepsilon)e^{-\nu s}.
\]
Then $|\sigma'| \le C_f|y|_* + C_L\varepsilon e^{-\mu s} \le C_3(\delta + \varepsilon)e^{-\nu s}$, so $|\sigma(s) - \sigma(0)| \le C_3(\delta + \varepsilon)/\nu$. If $\delta + \varepsilon$ is so small that $C_0\delta + C_3(\delta + \varepsilon)/\nu < \sigma_2/4$ and $C_0\delta + C_2(\delta + \varepsilon) < r_2/4$, then $\zeta(s)$ stays in the image of $\R \times [-\sigma_2/4, \sigma_2/4] \times \{|y|_* \le r_2/4\}$ on $[0, S_1)$, a compact set inside the image of $\R \times (-\sigma_2/2, \sigma_2/2) \times \{|y|_* < r_2/2\}$, and by continuity $S_1 = S$. The two distance bounds follow, since $|\Phi(\phi, \sigma, y) - e^{i\phi}\zeta^*| \le |\xi(\sigma) - \zeta^*| + |Qy|$ and $|\Phi(\phi, \sigma, y) - e^{i\phi}\xi(\sigma)| = |Qy|$. If $S = \infty$, $\sigma(s)$ converges to some $\sigma_\infty$ with $|\sigma(s) - \sigma_\infty| \le C_3(\delta + \varepsilon)e^{-\nu s}/\nu$, and $|\zeta(s) - e^{i\phi(s)}\xi(\sigma_\infty)| \le |\xi(\sigma(s)) - \xi(\sigma_\infty)| + |Qy(s)|$.
\end{proof}

\begin{proof}[Proof of Theorem~\ref{thm:nonlinear}]
Apply Proposition~\ref{prop:forced} with $\varepsilon = 0$, $\mu = 2$ and $\nu = 1/2$ to the solution on its maximal interval of existence. It stays in a compact set without collisions, so it exists for all $s \ge 0$, and $|\sigma(s) - \sigma_\infty| + |y(s)| \le C\delta e^{-s/2}$. For the sharp rate write $y' = B(\sigma_\infty)y + r(s)$ with $r = (B(\sigma(s)) - B(\sigma_\infty))y + (g(\sigma, y) - B(\sigma)y)$, so $|r(s)| \le C\delta e^{-s/2}|y(s)|$. By Lemma~\ref{lem:coords} and the variation of constants formula, $e^{s}|y(s)| \le K_B|y(0)| + K_BC\delta\int_0^se^{-\tau/2}e^{\tau}|y(\tau)|\,d\tau$, and Gronwall's inequality gives $|y(s)| \le K_B|y(0)|e^{2K_BC\delta}e^{-s} \le C\delta e^{-s}$. Then $|\sigma(s) - \sigma_\infty| \le \int_s^\infty C_f|y| \le C\delta e^{-s}$. With $\beta_\infty = b(\sigma_\infty) - b^*$, $\phi' = h(\sigma, y) = -\beta_\infty + O(|\sigma - \sigma_\infty| + |y|) = -\beta_\infty + O(\delta e^{-s})$, so $\phi(s) + \beta_\infty s$ converges to some $\phi_\infty$ at the rate $\delta e^{-s}$, and the estimate follows from $\zeta(s) = e^{i\phi(s)}(\xi(\sigma(s)) + Qy(s))$.

For the energy: the sum $\sum_{j<k}\Gamma_j\Gamma_k$ vanishes exactly for the circulations of the theorems, so $H(\lambda\zeta) = H(\zeta)$ for $\lambda > 0$; with the invariance of $H$ under rotations and its conservation by \eqref{eq:bs}, $H(\zeta(s))$ is constant along \eqref{eq:exp}. As $e^{-i\phi(s)}\zeta(s) \to \xi(\sigma_\infty)$, $H(\xi(\sigma_\infty)) = H(\zeta(0))$. For $\delta_0$ small, $|\sigma_\infty| \le C\delta_0 \le \sigma_0$, and since $\sigma_0$ was chosen so that $H(\xi(\sigma))$ is strictly monotone on $|\sigma| \le \sigma_0$, $\sigma_\infty$ is determined by $H(\zeta(0))$.
\end{proof}

\begin{remark}[Why the energy varies along the family]\label{rem:energy}
Differentiating the identity $\nabla H(\zeta)\cdot E(\zeta) = 0$, which expresses the conservation of $H$ and its invariance under rotations and scalings, at a zero of $E$ shows that $\nabla H(\zeta^*)$ is a left null vector of $DE$; it is not zero, since $2\pi V$ is, up to sign, the Hamiltonian vector field of $H$ for the nondegenerate form $\sum_j\Gamma_j\,dx_j \wedge dy_j$ (proof of Lemma~\ref{lem:spectrum}) and $2\pi V(\zeta^*) = -(1 - ib^*)\zeta^* \ne 0$. When $b'(0) \ne 0$, the kernel of $DE$ is spanned by $i\zeta^*$ alone, and the range of $DE$ is $W_1 \oplus \R\,i\zeta^*$, a hyperplane that does not contain $\xi'(0)$. As $\nabla H(\zeta^*)$ vanishes on that range, $dH/d\sigma = \nabla H(\zeta^*)\cdot\xi'(0) \ne 0$. The program encloses $b'(0)$ away from $0$ as well ($0.94234865\ldots$ and $-9.6795840\ldots$) and checks that $\nabla H(\zeta^*)^{\mathsf T}DE$ contains $0$. The limit of a perturbed expansion is therefore fixed by a conserved quantity, which Section~\ref{sec:numerics} illustrates.
\end{remark}

\section{Vortex patches}\label{sec:patches}

We now replace the point vortices by patches of vorticity. The Euler equations $\partial_t\omega + u\cdot\nabla\omega = 0$, $u = K*\omega$, $K(x) = x^\perp/(2\pi|x|^2)$, move point vortices by \eqref{eq:bs}, and have a unique global solution for bounded, compactly supported initial vorticity (Yudovich); a solution $\omega(t) = \sum_j\omega_j(t)$ is the sum of the parts transported from the initial pieces $\omega_j$. Let $\zeta^*$, $b^*$ and $\Gamma$ be as in Section~\ref{sec:nonlinear}, and put $y = \zeta^*/\sqrt\pi$. With the origin of time moved to the point of emergence, $\rho^2 = t/\pi$, the expansion of Corollary~\ref{cor:physical} is $x^*(t) = \sqrt t\,e^{-i(b^*/2)\ln(t/\pi)}y$, a motion of point vortices with circulations $-\Gamma$. Let $V = \{z \in \C^N: \sum_j\Gamma_jz_j = 0\}$, which contains $y$ and $iy$, and let $S = \{z \in V: \re\langle z - y, y\rangle = \re\langle z - y, iy\rangle = 0\}$ with $\langle z, w\rangle = \sum_j\overline{z_j}w_j$, a slice of dimension $2N - 4$ through $y$, transverse to rotations and dilations.

\begin{theorem}[Confinement of vortex patches]\label{thm:patches}
Let $\zeta^*$ be the configuration of Theorem~\ref{thm:four} or of Theorem~\ref{thm:five}, with $y$ and $S$ as above. Let $\varepsilon > 0$ be small enough, $M > 0$, $r > 0$ and $\nu \in (0, 1)$. There is $T$ such that for every $t_0 > T$ the following holds. Let $\omega_1, \ldots, \omega_N$ be bounded functions of definite sign with $\operatorname{supp}\omega_j \subseteq B(\sqrt{t_0}\,y_j, r)$, $\|\omega_j\|_{L^{2+\varepsilon}} \le M$, $\int\omega_j = -\Gamma_j$ and $\sum_j\int x\,\omega_j(x)\,dx = 0$. Let $\omega(t) = \sum_j\omega_j(t)$, $t \ge t_0$, be the solution of the Euler equations with $\omega(t_0) = \sum_j\omega_j$, and let $\bar x_j(t) = -\Gamma_j^{-1}\int x\,\omega_j(t, x)\,dx$ be the centres of vorticity. Then for all $t \ge t_0$ and all $j$:
\begin{enumerate}
\item[(i)] $\omega_j(t)$ has the sign of $-\Gamma_j$ and $\operatorname{supp}\omega_j(t) \subseteq B(\bar x_j(t), t^{1/4 + \varepsilon})$;
\item[(ii)] $\bar x(t) = e^{i\beta(t)}\gamma(t)z(t)$ with uniquely determined $\beta(t) \in \R/2\pi\mathbb Z$, $\gamma(t) > 0$ and $z(t) \in S$, and $|z(t) - y| < \varepsilon$, $|\gamma(t)/\sqrt t - 1| < \varepsilon$;
\item[(iii)] there are $\sigma_\infty$ with $|\sigma_\infty| < \varepsilon$, a continuous function $\phi$ and a constant $C$ with $\bigl|\bar x(t)/\sqrt t - e^{i\phi(t)}\xi(\sigma_\infty)/\sqrt\pi\bigr| \le Ct^{-\nu/2}$.
\end{enumerate}
\end{theorem}

Parts (i) and (ii) are the conclusion of Zbarsky's theorem for three vortices \cite[Theorem~2]{zbarsky2021}, with his generic expanding triangle replaced by the configurations of Theorems~\ref{thm:four} and~\ref{thm:five}. Part (iii) adds that the centres of vorticity move asymptotically like the self-similar expansion of a single member of the family. Dávila, del Pino, Musso and Parmeshwar \cite{dpmp2024} construct smooth solutions whose vorticity stays in balls of fixed radius around an expanding three-vortex spiral, and remark that their construction applies to four and five vortices under analogous conditions on the circulations. That is a construction of particular solutions; Theorem~\ref{thm:patches} is about every initial datum of the stated kind.

\begin{proof}
We follow the proof of \cite[Theorem~2]{zbarsky2021}, with the equation numbers of arXiv:1912.10862v2, and indicate every change. Zbarsky measures time in units of $2\pi$, so that $K(x) = x^\perp/|x|^2$; in his units the reference configuration is $\sqrt{2\pi}\,y$ and the slice $\sqrt{2\pi}\,S$, and (i) and (ii) in his units imply them in ours, since $t/2\pi < t$. We run his argument in his units, except in step (b), where the conversion changes the bounds only by the factor $\sqrt{2\pi}$, absorbed by enlarging $T$. His bootstrap assumptions, for $j = 1, \ldots, N$, are:
(B1) $\bar x = e^{i\beta}\gamma z$ with $z \in S$, $|z - y| < \varepsilon^2$ and $|\gamma/\sqrt t - 1| < \varepsilon$;
(B2) $I_{2,j} < t^{\varepsilon/2}$;
(B3) $I_{k,j} < t^{k(1 + \varepsilon)/4}$ for a fixed large even $k$;
(B4) $\omega_j$ has definite sign and $\|\omega_j\|_{L^{2+\varepsilon}} \le M$;
(B5) $\operatorname{supp}\omega_j \subseteq B(\bar x_j, t^{1/4+\varepsilon})$.
Here $I_{k,j} = \int|x - \bar x_j|^k|\omega_j(x)|\,dx$. They hold at $t_0$ when $T$ is large; for (B1) this follows from the slice decomposition below, since $|\bar x(t_0) - \sqrt{t_0}\,y| \le \sqrt N\,r$. Suppose that one of them fails at a first time $T_* < \infty$. On $[t_0, T_*]$ all of them hold.

\emph{Where $N = 3$ entered.} Zbarsky's proof uses three vortices in one place only: to recover (B1), at the end of his Section~3. There he uses that the angular impulse and the energy of the centres are nearly conserved, and that for three vortices these two quantities, with the rotation and the dilation, are coordinates near $y$ \cite[Lemma~1]{zbarsky2021}. For $N \ge 4$ the configurations modulo rotation and dilation have dimension $2N - 4 > 2$, and two nearly conserved quantities cannot pin them down. We use Proposition~\ref{prop:forced} instead.

\emph{(a) The centres move as point vortices up to a small error.} By (B1), $c\sqrt t \le |\bar x_i - \bar x_j| \le C\sqrt t$ and $|\bar x_j| \le C\sqrt t$ for $i \ne j$, with $c$, $C$ depending only on $y$. Zbarsky's computation on his pp.~7--8, with the other patches summed over $j \ne i$, expands the kernel between two patches about their centres of vorticity. The first-order terms vanish because $\bar x_j$ are the centres, the self-interaction vanishes by antisymmetry, and the remainder is $O((I_{2,i} + I_{2,j})t^{-3/2})$. Hence
\[
  \frac{d\bar x}{dt} = V_{-\Gamma}(\bar x) + e(t),\qquad |e(t)| \le Kt^{\varepsilon/2 - 3/2}\quad(t_0 \le t \le T_*),
\]
with $V_{-\Gamma}$ the velocity \eqref{eq:bs} for the circulations $-\Gamma$ (in the units of \eqref{eq:bs}) and $K$ depending on $y$, $\varepsilon$, $M$ and $r$ but not on $T_*$. Also $|d\bar x_j/dt| = O(t^{-1/2})$.

\emph{(b) Recovering (B1).} Put $\zeta(s) = \rho^{-1}e^{ib^*s}\bar x(t)$ with $\rho = \sqrt{t/\pi}$ and $s = \ln\rho$. The computation after Proposition~\ref{prop:forced} uses only $d\rho/dt = 1/(2\pi\rho)$, so $\zeta' = -E(\zeta) + 2\pi\rho e^{ib^*s}e(t)$, and $2\pi\rho|e(t)| \le 2\pi K\rho(\pi\rho^2)^{\varepsilon/2 - 3/2} \le 2\pi K e^{-(2 - \varepsilon)s}$. In the shifted time $\tau = s - s_0$, $s_0 = \ln\rho(t_0)$, the hypothesis of Proposition~\ref{prop:forced} holds with $\mu = 2 - \varepsilon$, forcing $\varepsilon_1 = 2\pi K\rho(t_0)^{\varepsilon - 2}$, and $\delta = \operatorname{dist}(\zeta(s_0), \mathcal O) \le |\bar x(t_0) - \sqrt{t_0}\,y|/\rho(t_0) \le \sqrt{\pi N}\,r\,t_0^{-1/2}$. Both tend to $0$ as $t_0 \to \infty$, so for $T$ large the proposition, applied with the given $\nu$ on $[s_0, \ln\rho(T_*))$ and extended to $T_*$ by continuity, gives $\operatorname{dist}(\zeta(s), \mathcal O) \le C_1(\delta + \varepsilon_1)$ on $[t_0, T_*]$. Now $\min_\phi|\bar x(t) - \sqrt t\,e^{i\phi}y| = \sqrt{t/\pi}\,\operatorname{dist}(\zeta(s), \mathcal O)$.

The slice decomposition goes as follows. The map $(\beta, \gamma, z) \mapsto e^{i\beta}\gamma z$ from $\R \times (0, \infty) \times S$ to $V$ has an invertible differential at $(0, 1, y)$, because its image is $\R\,iy \oplus \R\,y \oplus T_yS = V$. By the inverse function theorem, equivariance under rotations and dilations, and the argument for injectivity in the proof of Lemma~\ref{lem:coords}, there are $\eta$, $C_S > 0$ with the following property. Every $w \in V$ with $|w - \lambda e^{i\phi}y| \le \eta\lambda$ for some $\lambda > 0$ and $\phi$ is $e^{i\beta}\gamma z$ with unique $\beta$ near $\phi$ (mod $2\pi$), $\gamma > 0$ and $z \in S$, and $|z - y| + |\gamma/\lambda - 1| \le C_S|w - \lambda e^{i\phi}y|/\lambda$. Here $\bar x(t) \in V$, because $\sum_j\Gamma_j\bar x_j = -\sum_j\int x\,\omega_j$ is conserved and vanishes at $t_0$. Hence $|z - y| + |\gamma/\sqrt t - 1| \le C_SC_1(\delta + \varepsilon_1)/\sqrt\pi$ on $[t_0, T_*]$, which is less than $\varepsilon^2/(2\sqrt{2\pi})$ when $T$ is large. So (B1) holds on $[t_0, T_*]$ with strict inequalities in either unit of time, in particular at $T_*$.

\emph{(c) The other assumptions.} Zbarsky's Sections~4, 4.1 and~5 treat one patch at a time and use (B1) only through (a). For $N$ patches, $v_2 + v_3$ becomes $\sum_{j\ne1}v_j$ in (3)--(5), (19), (34) and (35). The renormalizing functions $f_{k,2}$, $f_{k,3}$ in (11)--(18) become $f_{k,j}$, one for each $j \ne 1$, and the angles $\theta_2$, $\theta_3$ in (29)--(31) become $\theta_j$, $j \ne 1$. For each $j$ separately, the principal term in (30) averages to zero over one turn of $p$ about $\bar x_1$. These sums have $N - 1$ terms, which changes only constants. The concentration estimate (33) uses the conserved energy $\iint\log|x - y|\,\omega(x)\omega(y)\,dx\,dy$. There the cross terms contribute $\sum_{i\ne j}\Gamma_i\Gamma_j(\ln t)/2 + O(1) = O(1)$, because $\sum_{i<j}\Gamma_i\Gamma_j = 0$ for our circulations, which the program checks in exact arithmetic. One term is missing from (34) for any $N$: $f_{k,j}$ depends on $t$ through $\bar x_1$ and $\bar x_j$, and the omitted $\int\partial_tf_{k,j}\,\omega_1$ is $O((t^{-1/4+\varepsilon} + t^{-1/2+2\varepsilon})I_{k,1})$ by (B5), $|d\bar x_1/dt| = O(t^{-1/2})$ and (29). That is absorbed in the error $O(\delta I_{k,1})$ of (12) for large $t$. With these changes his arguments recover (B2), (B3) and (B5) at $T_*$, in both the case $T_* \ge t_0 + t_0^{9/10}$ (his Section~4) and the case $T_* < t_0 + t_0^{9/10}$ (his Section~5). (B4) holds because the vorticity is transported by a divergence-free field. This contradicts the choice of $T_*$, so every assumption holds for all $t \ge t_0$, which gives (i) and (ii).

\emph{(d) Part (iii).} Now the estimate of (a) holds for all $t \ge t_0$, so Proposition~\ref{prop:forced} applies with $S = \infty$. It gives $\sigma_\infty$ with $|\sigma_\infty| \le C_1(\delta + \varepsilon_1) < \varepsilon$ and $\phi_1$ with $|\zeta(s) - e^{i\phi_1(s)}\xi(\sigma_\infty)| \le C_1(\delta + \varepsilon_1)e^{-\nu(s - s_0)}$. Since $\bar x(t)/\sqrt t = e^{-ib^*s}\zeta(s)/\sqrt\pi$ and $e^{-\nu(s - s_0)} = (t_0/t)^{\nu/2}$, (iii) follows with $\phi = \phi_1 - b^*s$.
\end{proof}

\section{Numerical illustration, controls and a sample}\label{sec:numerics}

\emph{Controls.} The same program certifies two controls, with the same side conditions. A four-vortex collapse with circulations $(1, -1003/1000, 13/10, \Gamma_4)$ has $c = -47.7346\ldots < 0$: its remaining pair is real, $k < 0$, and its reversal is unstable. A three-vortex collapse has $\tr((DE - I)^2) = 4 - 2b^2$ to within its enclosure, the contribution of the six forced eigenvalues alone, as Lemma~\ref{lem:spectrum} requires for $N = 3$.

\emph{Direct integration} (binary64, DOP853, relative tolerance $10^{-12}$, and $10^{-13}$ for the longest run; not part of the proof). The four-vortex expansion of Corollary~\ref{cor:expansion}, started from the midpoint of the certified box, keeps its shape to $4\cdot10^{-15}$ while its size grows $178$-fold; a random perturbation of $10^{-5}$ in each coordinate (a shape deviation of $3.3\cdot10^{-5}$) stays below $2\cdot10^{-4}$, settling on a nearby member of the family, the neutral direction. Measured against the member with the same energy, found by least squares as Theorem~\ref{thm:nonlinear} prescribes, a perturbation of $10^{-4}$ in each coordinate has a shape deviation (modulo translation, rotation and scaling, so weaker than the quantity in the theorem) that decays like the inverse of the size. Followed from $5$-fold to $5\cdot10^8$-fold growth, deviation times size stays between $3.0\cdot10^{-4}$ and $6.3\cdot10^{-4}$ for four vortices, and between $3.1\cdot10^{-3}$ and $2.1\cdot10^{-2}$ for five, where two oscillating modes of nearly equal frequency beat with a period of about $24$ in $\ln\rho$. The deviation from $\zeta^*$ itself levels off, near $1.1\cdot10^{-3}$ and $5.8\cdot10^{-3}$, and $H$ is conserved to $4\cdot10^{-14}$. The reversal of the unstable control, perturbed by $10^{-8}$ in each coordinate, departs from its shape by $1.5\cdot10^{-3}$ when its size has grown eightfold and by $0.3$ at sixteenfold. A self-similar collapse, by contrast, departs from self-similarity in binary64 under round-off alone \cite{lk2015, cwr2026}.

\emph{How common.} A least-squares solver from random starts, in binary64 with fixed seeds, gave $342$ converged four-vortex collapses (not checked for duplicates), of which $52$ have $c > 0$, and $543$ five-vortex collapses, of which $32$ have every remaining exponent with positive real part (\texttt{survey\_\allowbreak expansions.py}). This is a sample of what a naive search finds, not a measure on the collapse set. Of the $52$ four-vortex cases, $37$ have the remaining pair on $\re k = 1$, as in Theorem~\ref{thm:four}, and $15$ a real pair $0 < k < 2$. Of the $32$ five-vortex cases, $10$ have both remaining pairs on $\re k = 1$, as in Theorem~\ref{thm:five}, $16$ a quadruple $1 \pm x \pm iy$ with $0 < x < 1$, and $6$ a real pair $0 < k < 2$.

\begin{table}[t]
\centering
\small
\begin{tabular}{@{}>{\raggedright\arraybackslash}p{0.3\textwidth}>{\raggedright\arraybackslash}p{0.55\textwidth}r@{}}
\toprule
program & what it checks & checks \\
\midrule
\texttt{verify\_\allowbreak stable\_\allowbreak expansion.py} & Theorems~\ref{thm:four}, \ref{thm:five}, the hypotheses of Lemma~\ref{lem:spectrum} and of Theorem~\ref{thm:nonlinear} on the boxes, controls, integrations & $58$ \\
\texttt{survey\_\allowbreak expansions.py} & the random sample of Section~\ref{sec:numerics} (numerical) & \\
\bottomrule
\end{tabular}
\caption{The programs. Each prints every check and stops with an error if one fails; the output is in \texttt{data/}.}\label{tab:programs}
\end{table}

\section{Discussion}\label{sec:discussion}

The theorems give self-similarly expanding configurations of four and five vortices that are linearly asymptotically stable modulo the rotation and the family of expansions with the same circulations, with a spectral gap. Theorem~\ref{thm:nonlinear} makes this nonlinear: in similarity variables, nearby motions converge, up to a rotation drifting at the rate $b(\sigma_\infty) - b^*$ and at the rate the linearization predicts, to the member of the family with their own energy, and in the original variables they stay within a bounded distance of an exactly self-similar expansion. Zbarsky writes that orbital linear stability, ``which may not be hard to obtain for some systems, is not enough'' \cite{zbarsky2021}. What his argument needs is an estimate for approximate solutions whose error is small and decays, with constants that do not depend on the length of time, and Proposition~\ref{prop:forced} is such an estimate. Theorem~\ref{thm:patches} is then his theorem for these two configurations. Its proof rests on his estimates for single patches, which we have followed and adapted but not rederived in full, so it should be read together with \cite{zbarsky2021}; the changes are listed in the proof, and one omission in his estimate (34), which does not affect his conclusion, is corrected there. An independent check by a specialist in vortex patches would be valuable before the result is relied on.

The examples are Euler vortices. The $\alpha$-models, where a vortex induces the velocity $\Gamma r^{-\alpha-1}/(2\pi)$, have the same structure after the obvious changes of homogeneity, and we have not treated them. Nor have we asked which fraction of the collapse set reverses into stable expansions, or whether stable expansions exist for every $N \ge 4$; the sample of Section~\ref{sec:numerics} suggests that they are common for $N = 4$ and $5$.

\medskip
{\raggedright\noindent\textbf{Data availability.} The programs \texttt{verify\_\allowbreak stable\_\allowbreak expansion.py} and \texttt{survey\_\allowbreak expansions.py}, with the modules \texttt{certify\_\allowbreak ball\_\allowbreak ad.py} and \texttt{certify\_\allowbreak pipeline.py} of \cite{mw2026}, their inputs and their output will be in the repository \url{https://github.com/ChaseHendrick/stable-expansion}: the programs in \texttt{code/} and their inputs and output in \texttt{data/}.\par}

\smallskip
\noindent\textbf{Funding.} This research received no external funding.

\smallskip
\noindent{\footnotesize This work was prepared with AI assistance. The author takes full responsibility for its content.}

\begin{thebibliography}{14}

\bibitem{aref2010} H. Aref, Self-similar motion of three point vortices, \emph{Phys. Fluids} \textbf{22} (2010) 057104.

\bibitem{aulbach1984} B. Aulbach, \emph{Continuous and Discrete Dynamics near Manifolds of Equilibria}, Lecture Notes in Mathematics \textbf{1058}, Springer, Berlin (1984); doi:10.1007/BFb0071569.

\bibitem{dpmp2024} J. D\'avila, M. del Pino, M. Musso and S. Parmeshwar, An expanding self-similar vortex configuration for the 2D Euler equations, preprint, arXiv:2410.18220 (2024).

\bibitem{gotoda2021} T. Gotoda, Self-similar motions and related relative equilibria in the $N$-point vortex system, \emph{J. Dyn. Differ. Equ.} \textbf{33} (2021) 1759--1777; arXiv:2002.09624.

\bibitem{cwr2026} C. Hendrick, Point-vortex collapse without rotation: a cluster mechanism, a phase diagram and a continuum limit, draft (2026), Sect.~5.3, at \url{https://github.com/ChaseHendrick/GENChase/tree/main/papers/collapse-without-rotation}.

\bibitem{mw2026} C. Hendrick, Minimal winding in the self-similar collapse of point vortices, preprint (2026); programs and data, release v2.0.0, doi:10.5281/zenodo.22963796, at \url{https://github.com/ChaseHendrick/minimal-winding}.

\bibitem{iy2023} T. Iwayama and T. Yajima, Linear stability of self-similar motions of three point vortices in a generalized two-dimensional fluid system, \emph{J. Phys. Soc. Jpn.} \textbf{92} (2023) 084401.

\bibitem{ks2022} S. S. Kallyadan and P. Shukla, Self-similar vortex configurations: collapse, expansion, and rigid-vortex motion, \emph{Phys. Rev. Fluids} \textbf{7} (2022) 114701.

\bibitem{kimura1987} Y. Kimura, Similarity solution of two-dimensional point vortices, \emph{J. Phys. Soc. Jpn.} \textbf{56} (1987) 2024--2030.

\bibitem{lk2015} M. Lewkowicz and H. Kudela, Collapse of $n$ vortices, preprint, arXiv:1512.05116 (2015).

\bibitem{ns1979} E. A. Novikov and Yu. B. Sedov, Vortex collapse, \emph{Sov. Phys. JETP} \textbf{50} (1979) 297--301 [\emph{Zh. Eksp. Teor. Fiz.} \textbf{77} (1979) 588--597].

\bibitem{tt1988} J. Tavantzis and L. Ting, The dynamics of three vortices revisited, \emph{Phys. Fluids} \textbf{31} (1988) 1392--1409.

\bibitem{zbarsky2021} S. Zbarsky, From point vortices to vortex patches in self-similar expanding configurations, \emph{Commun. Math. Phys.} \textbf{388} (2021) 707--733; arXiv:1912.10862.

\bibitem{zbarsky2024} S. Zbarsky, Bounds on growth and impossibility of collapse for point vortex systems, preprint, arXiv:2402.07316 (2024).

\end{thebibliography}

\end{document}
